
\clearpage
\VerdierPhysicalPageAnchor{0126}%
\begingroup
\small
\setlength{\emergencystretch}{1.5em}%
\setlength{\abovedisplayskip}{4pt plus 2pt minus 2pt}%
\setlength{\belowdisplayskip}{4pt plus 2pt minus 2pt}%

\VerdierSemanticAnchor{ch2-1-3-7}
\noindent\textbf{1.3.7.}\quad
\VerdierSemanticAnchor{ch2-1-3-7-multi-additive-functor}
$F\colon \prod_{i\in[n]} C_i\longrightarrow C$
を多加法的関手とする。$C$ において可算積が表現可能であると仮定する。
（第 I 章、3.4.3 および 3.4.4）から、関手 $F$ の複体のホモトピー圏への
$\Pi$-拡張

\VerdierSemanticAnchor{ch2-1-3-7-pi-extension}
\[
  K_{\Pi}F\colon \prod_{i\in[n]} K(C_i)\longrightarrow K(C),
\]

\VerdierSemanticAnchor{ch2-1-3-7-multi-exact}
\noindent は多重完全関手であることが直ちに従う。

\VerdierSemanticAnchor{ch2-1-3-7-koszul-twist}
この関手は符号規則（第 I 章、1.6.6）に従い、Koszul 余サイクルで捩られた
$\mathbb Z^n$-関手である。

\VerdierSemanticAnchor{ch2-1-3-7-sigma-extension}
関手 $K_{\Sigma}F$ についても同様の結果が成り立つ。

\VerdierSemanticAnchor{ch2-1-3-8}
\medskip
\noindent\textbf{1.3.8.}\quad
\VerdierSemanticAnchor{ch2-1-3-8-canonical-opposite-isomorphism}
$C$ を加法圏とし、$\iota\colon K(C)^\circ\longrightarrow K(C^\circ)$ を
標準同型 $\iota=\iota_1$（第 I 章、2.5.13）とする。

\VerdierSemanticAnchor{ch2-1-3-8-triangulated-opposite}
圏 $K(C)^\circ$ に (1.1.7) で述べた三角構造を入れる。

\VerdierSemanticAnchor{ch2-1-3-8-exactness}
（第 I 章、3.4.1）により、$\iota$ は三角圏の間の完全関手である。
したがって $\iota$ は三角圏の同型である。

\VerdierSemanticAnchor{ch2-1-3-8-identification}
以下、特に断らない限り、同型 $\iota$ によって三角圏 $K(C)^\circ$ と
$K(C^\circ)$ を同一視する。
\endgroup


